\chapter{Evaluating Synthetic Data}
\label{ch:evaluation}

\begin{tcolorbox}[feynbox={Sky}{BBg},title={Before and After}]
\textbf{Before this chapter you need:} permutations and independence
(\S\ref{sec:joint}), the CDF (\S\ref{sec:distributions}), hypothesis testing,
$p$-values, power and multiple testing
(\S\ref{sec:hypothesis}--\S\ref{sec:multipletesting}), autocorrelation and Hurst
(\S\ref{sec:acf}, \S\ref{sec:hurst}), tail index (\S\ref{sec:moments}),
Wasserstein distance (\S\ref{sec:wasserstein}), the stylised facts
(\S\ref{sec:stylisedfacts}).\\
\textbf{After it you will understand:} the single structural property ---
permutation invariance --- that determines what any metric can possibly detect,
why nine of nineteen ranked metrics shared it and handed victory to a shuffled
deck, and the taxonomy that repairs it.
\end{tcolorbox}

Training defines what a model optimises; evaluation defines what counts as
success. This chapter argues that the second is the harder problem, and that
getting it wrong is easy to do without noticing.

%% ─────────────────────────────────────────────────────────────────────────────
\section{What Evaluation Must Answer}
\label{sec:evalquestions}

\situation{
Someone hands you a synthetic dataset and says it looks like the real thing.
Looks how? Under what test? A picture of two overlapping histograms is not an
answer, and neither is a single number without an account of what it can miss.
}

\problem{
There is no likelihood to report (\S\ref{sec:mletraining}), so evaluation must
work entirely from samples. The danger is not that a metric gives a wrong number;
it is that a metric gives a right number to a question nobody meant to ask.
}

\workedarith{
\textbf{Three separable questions.}

\medskip
\begin{tabular}{@{}lll@{}}
\toprule
Question & What it asks & Family \\
\midrule
Are the values right? & marginal distribution of $r_t$ & fidelity \\
Is the ordering right? & dependence between $r_t$ and $r_{t-k}$ & temporal \\
Is it useful? & does a model fitted on it work on real data & downstream \\
\bottomrule
\end{tabular}

\medskip
\textbf{They are genuinely independent.} A series can pass any one and fail the
others:

\begin{itemize}
  \item Shuffled real data: perfect on the first, fails the second.
  \item Gaussian noise with the right variance: passes crude versions of the
        first, fails on tails and on all of the second.
  \item A memorised replay of the training set: passes all three and generates
        nothing new.
\end{itemize}

\medskip
\textbf{The consequence.} No single number answers all three, and averaging them
into one lets a model strong in a well-populated family outvote its weakness
elsewhere. Section~\ref{sec:permutation} shows this happening.
}

\formaldef{
An evaluation framework for synthetic data comprises:
\begin{itemize}
  \item \textbf{fidelity} --- agreement of marginal distributions;
  \item \textbf{temporal} --- agreement of dependence structure;
  \item \textbf{downstream utility} --- performance of a model trained on
        synthetic data and tested on real (Chapter~\ref{ch:downstream});
  \item \textbf{controls} --- reference series with known properties, against
        which the metrics themselves are validated.
\end{itemize}
The last is the one usually omitted, and it is what this chapter argues cannot be.
}

\analogybreak{
Evaluation frameworks are themselves models, with assumptions that can fail. Every
metric here assumes stationarity (\S\ref{sec:stationarity}) --- that the test
distribution matches the training one --- which the MOEX structural break of
\S\ref{sec:breaks} violates outright. A framework that never tests itself is in the
same position as a generator that is never evaluated.
}

\whyhere{
This is the methodological contribution of the thesis. The metrics themselves are
standard; the taxonomy that organises them, and the control that validated the
taxonomy, are what the framework adds.
}

%% ─────────────────────────────────────────────────────────────────────────────
\section{Distributional Similarity and Statistical Distances}
\label{sec:distances}

\situation{
Two bags of marbles. Same question asked three ways: is the largest colour
imbalance small, is the total imbalance small, or is the work needed to convert
one bag into the other small? Three sensible answers, three different numbers.
}

\problem{
``The distributions match'' needs a metric, and different metrics respond to
different kinds of disagreement. Choosing one without knowing what it is blind to
is how a framework acquires a hole.
}

\workedarith{
\textbf{Kolmogorov--Smirnov.} The largest vertical gap between two CDFs. Real
$\{1,3,5\}$, synthetic $\{2,3,4\}$:

\medskip
\begin{tabular}{@{}lrrr@{}}
\toprule
$x$ & $F_{\text{real}}$ & $F_{\text{syn}}$ & $|{\rm gap}|$ \\
\midrule
$1$ & $1/3$ & $0$   & $0.333$ \\
$2$ & $1/3$ & $1/3$ & $0$ \\
$3$ & $2/3$ & $2/3$ & $0$ \\
$4$ & $2/3$ & $1$   & $0.333$ \\
$5$ & $1$   & $1$   & $0$ \\
\bottomrule
\end{tabular}

$D_{\rm KS} = 0.333$, driven by the largest single discrepancy.

\medskip
\textbf{Wasserstein.} The average gap between sorted samples (\S\ref{sec:wasserstein}):
$W_1 = \tfrac13(1 + 0 + 1) = 0.667$. It accumulates every discrepancy rather than
taking the maximum.

\medskip
\textbf{Quantile MSE.} Squared differences at matched quantiles ---
like $W_1$ but penalising large gaps disproportionately.

\medskip
\textbf{Energy distance.} $2\E\norm{X-Y} - \E\norm{X-X'} - \E\norm{Y-Y'}$, zero
only when the distributions are identical, and sensitive to the whole shape.

\medskip
\textbf{Why KS is not used for ranking here.} It was designed for independent
data. Returns are dependent, so the effective sample size is far below the nominal
$n$ and KS $p$-values are badly anti-conservative. It is reported, not ranked.
}

\formaldef{
Fidelity metrics used in this project, all permutation-invariant:
\[
  \texttt{mean\_diff},\;
  \texttt{std\_diff},\;
  \texttt{wasserstein},\;
  \texttt{quantile\_mse},
\]
\[
  \texttt{tail\_index\_diff},\;
  \texttt{extreme\_events\_diff},\;
  \texttt{energy\_distance} .
\]
Seven metrics. Each depends only on the multiset of values, never on their order.
}

\analogybreak{
Every metric in this list has a floor set by the real data itself. Section
\ref{sec:wasserstein} showed the Wasserstein floor: a generator emitting the
constant $\mu$ scores $W_1$ equal to the real data's mean absolute deviation ---
$0.00770$ for BOVESPA --- and a genuine but imperfect generator can score worse.
Reading any of these metrics without knowing its floor invites exactly the wrong
conclusion.
}

\whyhere{
These seven constitute the fidelity family. \texttt{tail\_index\_diff} replaced
\texttt{kurtosis\_diff} for the reasons in \S\ref{sec:kurtosis}: the population
kurtosis is infinite at the measured tail indices, so ranking on its sample
estimate ranks on noise.
}

%% ─────────────────────────────────────────────────────────────────────────────
\section{Temporal Similarity}
\label{sec:temporalsim}

\situation{
Two pieces of music using exactly the same notes in the same proportions. One is a
melody; the other is the notes shuffled. Any measurement that counts notes calls
them identical.
}

\problem{
Fidelity metrics are blind to order by construction. Everything that makes a
return series a \emph{series} rather than a bag of numbers must be measured by a
second family, chosen so that shuffling changes the answer.
}

\workedarith{
\textbf{The test any temporal metric must pass.} Shuffle the real data. A temporal
metric must score the shuffle badly; a fidelity metric cannot.

\medskip
\begin{tabular}{@{}lrl@{}}
\toprule
Metric on shuffled real & Score & Correct? \\
\midrule
\texttt{wasserstein}       & $0.000\text{e}{+}00$ & yes, and useless \\
\texttt{kurtosis\_diff}    & $1.78\times10^{-15}$ & yes, and useless \\
\texttt{energy\_distance}  & $2.97\times10^{-5}$  & yes, and useless \\
\midrule
\texttt{acf\_returns\_mae}   & $0.043$ & penalised \\
\texttt{acf\_absolute\_mae}  & $0.076$ & penalised \\
\texttt{hurst\_diff}         & $0.052$ & penalised \\
\bottomrule
\end{tabular}

\medskip
\textbf{Read the top block.} Those are not failures of computation. The values are
exactly right: a permutation has the identical distribution, so the distributional
distance genuinely is zero. The metrics are working perfectly and answering a
question that cannot distinguish the two series.

\medskip
\textbf{Read the bottom block.} Shuffling destroys autocorrelation and long
memory, so these three respond. That responsiveness is the definition of a
temporal metric.
}

\formaldef{
Temporal metrics used here, all ordering-sensitive:
\[
  \texttt{acf\_returns\_mae},\;
  \texttt{acf\_absolute\_mae},\;
  \texttt{acf\_squared\_mae},\;
  \texttt{hurst\_diff},
\]
\[
  \texttt{resid\_kurtosis\_diff},\;
  \texttt{discriminative\_auc\_dist},\;
  \texttt{discriminative\_auc\_absz} .
\]
Seven metrics, matching the seven fidelity metrics in count --- which is what makes
an unweighted mean of the two families defensible.
}

\analogybreak{
Ordering-sensitivity is necessary but not sufficient. Three of these seven ---
the ACF triple --- measure closely related quantities, so the family is less
independent than its size suggests. Section~\ref{sec:multipletesting} made the
same point about Bonferroni: correlated tests do not count as separate evidence,
and seven correlated metrics are not seven independent checks.
}

\whyhere{
The pairing of facts 2 and 3 from \S\ref{sec:stylisedfacts} is why three ACF
metrics are needed rather than one. A generator producing independent noise
matches \texttt{acf\_returns\_mae} perfectly --- real returns are nearly
uncorrelated too --- while failing \texttt{acf\_absolute\_mae} completely. Only
computing both catches it.
}

%% ─────────────────────────────────────────────────────────────────────────────
\section{Discriminative and Predictive Scores}
\label{sec:discriminative}

\situation{
The practical test of a forgery is whether an expert can tell. Rather than
enumerating features, hand both to a classifier and see whether it can separate
them.
}

\problem{
Hand-chosen metrics test the properties we thought to test. A learned
discriminator searches for \emph{any} exploitable difference, including ones not
on the list --- which makes it the most general check available and the one most
easily misread.
}

\workedarith{
\textbf{The procedure.} Label real windows 1 and synthetic 0, train a classifier,
report its AUC on held-out data.

\medskip
\begin{tabular}{@{}lll@{}}
\toprule
AUC & Reading & \\
\midrule
$0.50$ & cannot separate & synthetic is indistinguishable \\
$0.75$ & separates well  & a detectable difference exists \\
$1.00$ & perfect         & trivially distinguishable \\
$0.30$ & anti-predictive & \emph{also} a detectable difference \\
\bottomrule
\end{tabular}

\medskip
\textbf{The direction bug.} Ranking ascending on raw AUC puts $0.30$ above
$0.50$ --- treating anti-predictive as better than indistinguishable. But a
classifier achieving $0.30$ can be made to achieve $0.70$ by inverting its output,
so it has found just as much signal. The fix is to rank on $|{\rm AUC} - 0.5|$.

\medskip
\textbf{Null calibration.} Two samples from the \emph{same} distribution do not
give exactly $0.5$. Fifteen real-versus-real half-splits gave
\[
  0.506 \pm 0.084 .
\]
An observed $0.62$ is therefore
\[
  z = \frac{0.62 - 0.506}{0.084} = 1.36 ,
\]
giving $p \approx 0.17$ --- not significant, despite looking well above $0.5$.

\medskip
\textbf{The CV shuffle bug.} Windows of 20 days overlap by 19 observations.
Shuffling the cross-validation folds places near-duplicate windows in both train
and test, leaking information. Measured effect on the null: $0.506 \to 0.584$.
Unshuffled cross-validation is required.
}

\formaldef{
The \textbf{discriminative score} is the AUC of a classifier trained to separate
real from synthetic windows. Two derived metrics are ranked:
\[
  \texttt{discriminative\_auc\_dist} = |{\rm AUC} - 0.5| ,
  \qquad
  \texttt{discriminative\_auc\_absz} = \left|\frac{{\rm AUC} - 0.506}{0.084}\right| .
\]
Both are minimised at the generator's optimum. The raw AUC is retained as
descriptive only.

The \textbf{predictive score} (TSTR in its sequence form) trains a forecaster on
synthetic data and measures its error on real data.
}

\analogybreak{
A discriminator that cannot separate the two proves only that \emph{this}
classifier, with these features and this much data, found nothing. A stronger
model might. AUC near the null is evidence of indistinguishability at a stated
level of scrutiny, never a proof of equality --- and \S\ref{sec:power} applies:
failure to detect is not evidence of absence when power is unknown.
}

\whyhere{
The three defects above --- direction, null calibration, and CV shuffling --- were
found in this project's own implementation, and all three would inflate apparent
performance if left uncorrected. They are documented as a caution: the
discriminative score is the most general metric in the framework and the easiest
to get quietly wrong.
}

%% ─────────────────────────────────────────────────────────────────────────────
\section{TSTR and TRTS}
\label{sec:tstr}

\situation{
A flight simulator is not judged on how much it resembles a cockpit. It is judged
on whether pilots trained in it can fly the aircraft.
}

\problem{
Distributional and temporal metrics ask whether synthetic data \emph{looks} right.
The question that matters to a user is whether it \emph{works} --- whether a model
built on it performs on reality.
}

\workedarith{
\textbf{The two protocols.}

\medskip
\begin{tabular}{@{}lll@{}}
\toprule
Protocol & Train on & Test on \\
\midrule
TSTR & synthetic & real \\
TRTS & real      & synthetic \\
TRTR & real      & real (the reference ceiling) \\
\bottomrule
\end{tabular}

\medskip
\textbf{TSTR is the one that matters.} It answers the practical question: if I had
only synthetic data, would the model I build be any good on the real world? TRTS
answers a narrower question --- whether the synthetic data is \emph{easy}, which a
degenerate generator satisfies trivially.

\medskip
\textbf{The reference.} TRTR gives QLIKE $\approx -8.134$ on this project's pooled
data. A synthetic dataset whose GARCH parameters transfer well should produce
TSTR QLIKE close to that. The gap is the cost of using synthetic data instead of
real.

\medskip
\textbf{Why the reference is essential.} Without TRTR, a TSTR QLIKE of $-7.9$ has
no interpretation. Against $-8.134$ it is a small, quantified shortfall.
}

\formaldef{
\textbf{TSTR} (Train on Synthetic, Test on Real) fits a downstream model on
synthetic data, applies it to real data, and scores it. \textbf{TRTS} reverses the
roles. \textbf{TRTR} trains and tests on real data and provides the upper
reference.

Here the downstream model is GARCH$(1,1)$: parameters are estimated on the
synthetic series, used to filter the real series, and the resulting conditional
VaR forecasts and QLIKE are scored --- the subject of
Chapter~\ref{ch:downstream}.
}

\analogybreak{
TSTR inherits every limitation of the downstream model it uses. Scoring through
GARCH tests whether the synthetic data supports \emph{GARCH-shaped} volatility
structure, and a generator producing a different but equally valid structure would
be penalised. The protocol measures usefulness for one purpose, and generalises no
further than that purpose.
}

\whyhere{
TSTR is what makes the evaluation more than a similarity exercise. A generator
matching every stylised fact and failing TSTR would be an interesting object and a
useless tool. QLIKE is the primary downstream signal because it is a proper
scoring rule and, at pooled $n \approx 2{,}480$, has adequate power ---
which at per-market $n \approx 126$ it demonstrably does not
(\S\ref{sec:qlike}).
}

%% ─────────────────────────────────────────────────────────────────────────────
\section{Permutation Invariance and the Shuffled Control}
\label{sec:permutation}

\situation{
You buy a shuffling machine. To test it, you check that the deck still holds 52
cards with four of each value --- and it does, every time. That check can never
fail, because shuffling does not remove cards. You have verified nothing about
shuffling.
}

\problem{
This is not an analogy. It is what the evaluation framework was doing, and the
control that exposed it is the central methodological finding of the thesis.
}

\workedarith{
\textbf{The control.} Take the pooled real test data and permute it:
\begin{quote}
\texttt{rng = numpy.random.default\_rng(42)}\\
\texttt{shuffled = rng.permutation(pooled\_real\_array).copy()}
\end{quote}
The result has \emph{identical} values in a random order. It is the best possible
fidelity generator --- it \emph{is} the real data --- and the worst possible
temporal one.

\medskip
\textbf{What it scored.}

\medskip
\begin{tabular}{@{}lr@{}}
\toprule
Metric & Shuffled control \\
\midrule
\texttt{wasserstein}        & $0.000\text{e}{+}00$ (exact zero) \\
\texttt{kurtosis\_diff}     & $1.78\times10^{-15}$ (floating-point noise) \\
\texttt{energy\_distance}   & $2.97\times10^{-5}$ \\
\texttt{mean\_diff}         & $< 10^{-14}$ \\
\texttt{quantile\_mse}      & $< 10^{-14}$ \\
\midrule
\texttt{acf\_absolute\_mae} & $0.076$ (correctly penalised) \\
\texttt{hurst\_diff}        & $0.052$ (correctly penalised) \\
\bottomrule
\end{tabular}

\medskip
\textbf{The verdict under the old design.} With nineteen metrics averaged
unweighted, the shuffled control scored an average rank of
\[
  1.24 ,
\]
against $2.47$ for a genuine generator. \textbf{The shuffled deck won the
benchmark.}

\medskip
\textbf{Why: count the invariant metrics.} Of the nineteen ranked metrics,
\emph{nine} were permutation-invariant --- the seven current fidelity metrics plus
\texttt{kurtosis\_diff} and \texttt{skewness\_diff}. Nine of nineteen, very nearly
half, scored the shuffle perfectly by construction. No weighting of the remaining
ten could overcome a near-perfect score on nine.

\medskip
\textbf{The arithmetic of the fix.} Split into two families of seven and average
the \emph{family ranks} rather than the metric ranks:
\[
  \texttt{composite\_rank} = \tfrac12\bigl(\texttt{fidelity\_rank} + \texttt{temporal\_rank}\bigr) .
\]
Now a perfect fidelity score buys at most half the composite, and the shuffled
control's temporal failure cannot be outvoted.
}

\formaldef{
A metric $m$ is \textbf{permutation-invariant} if
\[
  m\bigl(\pi(r_1,\ldots,r_n)\bigr) = m(r_1,\ldots,r_n)
\]
for every permutation $\pi$. Any metric depending only on the sorted values ---
the empirical CDF, moments, quantiles, order statistics --- is invariant. This
includes Wasserstein, energy distance, KS, kurtosis, skewness, mean, variance,
tail index and extreme-event rate.

A metric is \textbf{ordering-sensitive} if some permutation changes its value:
ACF, Hurst, GARCH-based quantities, and any classifier using time-window features.
}

\analogybreak{
The shuffled control is not a bad generator. It is the \emph{best possible}
fidelity generator, since it is the real data. Its failure on temporal metrics is
not a bug to be explained away --- it is the point. And it is not a rival to be
beaten: it is an instrument for testing the metrics, playing the role of a
positive control in an experiment.
}

\whyhere{
This finding restructured the evaluation. The control is now a permanent row in
every results table with NaN composite rank, and any future metric added to the
framework must be classified as invariant or ordering-sensitive before it can be
ranked. An evaluation without such a control cannot claim to have tested temporal
structure --- it can only claim not to have checked.
}

%% ─────────────────────────────────────────────────────────────────────────────
\section{The Metric Taxonomy}
\label{sec:taxonomy}

\situation{
A hospital reports outcomes. Mixing survival rates, waiting times and satisfaction
scores into one average produces a number that moves for reasons nobody can
identify. Grouping them first makes the number interpretable.
}

\problem{
The previous section established that families must be separated. This section
states the resulting taxonomy and, equally importantly, why seven metrics are
computed and deliberately excluded from ranking.
}

\workedarith{
\textbf{Fidelity, 7 metrics}, permutation-invariant:
\texttt{mean\_diff}, \texttt{std\_diff}, \texttt{wasserstein},
\texttt{quantile\_mse}, \texttt{tail\_index\_diff},
\texttt{extreme\_events\_diff}, \texttt{energy\_distance}.

\medskip
\textbf{Temporal, 7 metrics}, ordering-sensitive:
\texttt{acf\_returns\_mae}, \texttt{acf\_absolute\_mae},
\texttt{acf\_squared\_mae}, \texttt{hurst\_diff},
\texttt{resid\_kurtosis\_diff}, \texttt{discriminative\_auc\_dist},
\texttt{discriminative\_auc\_absz}.

\medskip
\textbf{Descriptive, 7 metrics}, computed and displayed, never ranked --- each for
a stated reason:

\medskip
\begin{tabular}{@{}ll@{}}
\toprule
Metric & Why excluded \\
\midrule
\texttt{kurtosis\_diff}      & $\E[X^4] = \infty$ at $\hat\alpha < 4$ \\
\texttt{skewness\_diff}      & $\E[X^3] = \infty$ in 4 of 5 markets \\
\texttt{arch\_pvalue\_diff}  & saturates: $|\Delta p| = 0.000000$ on real data \\
\texttt{arch\_stat\_diff}    & null sd $46.4$; too noisy to rank \\
\texttt{var\_coverage\_error}& Gaussian noise and shuffled real both score $0.0227$ \\
\texttt{garch\_persistence\_diff} & sd $0.31$ across seeds; range $0.000$--$0.985$ \\
\texttt{discriminative\_auc\_raw} & direction bug; superseded by dist and absz \\
\bottomrule
\end{tabular}

\medskip
\textbf{Reading the exclusions.} These are not metrics that performed badly. They
are metrics that cannot discriminate --- either because the quantity they estimate
does not exist, or because their noise exceeds their signal, or because a
degenerate control matches a genuine model on them. Computing and displaying them
while refusing to rank on them is more honest than either ranking on them or
dropping them silently.

\medskip
\textbf{The composite.}
\[
  \texttt{composite\_rank} = \tfrac12\bigl(\texttt{fidelity\_rank} + \texttt{temporal\_rank}\bigr) .
\]
Equal weight to both families regardless of how many metrics each contains.
\texttt{avg\_rank}, the legacy unweighted mean, is retained for comparability only
--- it is the score the shuffled control wins.
}

\formaldef{
Ranking uses \texttt{na\_option='bottom'}: a metric that fails to compute yields
NaN and ranks \emph{last}, so a failure counts as the worst outcome rather than
being silently excluded. This is what prevents the Hill $\alpha > 20$ rejection of
\S\ref{sec:hillfail} from becoming an advantage.

The shuffled control appears in every table with NaN composite rank --- present as
a reference, excluded from the competition.
}

\analogybreak{
Two families of seven is a defensible balance, not a derived optimum. The metrics
within each family are correlated (\S\ref{sec:temporalsim}), so neither family
represents seven independent pieces of evidence, and a different reasonable
taxonomy would give slightly different rankings. What the design guarantees is
narrower and more important: no amount of fidelity performance can conceal a
temporal failure.
}

\whyhere{
\texttt{composite\_rank} is the headline metric for model selection throughout the
thesis, and the pipeline selects on it rather than on \texttt{avg\_rank}. The
comment recording why sits in the code beside the selection, because this is
precisely the kind of decision that gets silently reverted during a refactor.
}

%% ─────────────────────────────────────────────────────────────────────────────
\section{Discriminative AUC: Null Calibration in Detail}
\label{sec:aucnull}

\situation{
A metal detector that beeps occasionally on clean sand. Before trusting it on a
beach, you must know its false-alarm rate on ground you know to be empty.
}

\problem{
Section~\ref{sec:discriminative} introduced the AUC null of
$0.506 \pm 0.084$. That number does real work in the framework, and where it comes
from --- and what it costs to get wrong --- deserves setting out fully.
}

\workedarith{
\textbf{How the null was measured.} Split the real data in half, label one half 1
and the other 0, and train the classifier. There is genuinely nothing to find, so
any AUC above $0.5$ is noise. Repeating fifteen times:
\[
  \text{AUC}_{\text{null}} = 0.506 \pm 0.084 .
\]

\medskip
\textbf{Why not exactly $0.5$.} Finite samples, and a classifier flexible enough to
fit some of the noise. The $+0.006$ offset is small; the $\pm 0.084$ spread is
not.

\medskip
\textbf{What the spread implies.} A $95\%$ interval is roughly
$0.506 \pm 1.96(0.084) = [0.34, 0.67]$. Any observed AUC inside that band is
consistent with the synthetic data being indistinguishable from real.

\medskip
\begin{tabular}{@{}lrrl@{}}
\toprule
Observed AUC & $z$ & $p$ & Verdict \\
\midrule
$0.55$ & $0.52$ & $0.60$ & indistinguishable \\
$0.62$ & $1.36$ & $0.17$ & indistinguishable \\
$0.70$ & $2.31$ & $0.02$ & detectably different \\
$0.90$ & $4.69$ & $<0.001$ & clearly different \\
\bottomrule
\end{tabular}

\medskip
\textbf{The cost of skipping this.} Judged against a theoretical $0.5$, an AUC of
$0.62$ looks like a clear failure. Judged against the measured null, it is
ordinary sampling noise. Without the calibration, the framework would report
differences that are not there.

\medskip
\textbf{And the leakage figure again.} Shuffling the cross-validation folds moves
the null from $0.506$ to $0.584$ --- a shift comparable to the entire effect one
might hope to measure. Overlapping windows make unshuffled CV mandatory here.
}

\formaldef{
\textbf{AUC} is $\Prob\bigl(\hat{p}(x_{\text{real}}) > \hat{p}(x_{\text{syn}})\bigr)$
for a randomly chosen pair. The generator's optimum is $0.5$.

The empirical null distribution --- $0.506 \pm 0.084$ from fifteen real-versus-real
splits with unshuffled $K$-fold CV --- is what
\texttt{discriminative\_auc\_absz} standardises against.
}

\analogybreak{
The null is specific to this classifier, these features, this window length and
this sample size. Change any of them and it must be re-measured. It is a property
of the measurement apparatus, not a universal constant, and quoting
$0.506 \pm 0.084$ in a different setting would be an error of the same kind as
quoting a Hurst exponent without naming the estimator (\S\ref{sec:roughvol}).
}

\whyhere{
Three defects were found and corrected in this metric: the direction bug,
the missing null calibration, and CV shuffling. All three inflate apparent
performance. They are reported in full because they are ordinary pitfalls rather
than exotic ones, and a framework that documents its own near-misses is more
trustworthy than one that presents only its final state.
}
